\documentclass[11pt]{article}
\usepackage[utf8]{inputenc}
\usepackage[T1]{fontenc}
\usepackage{amsmath,amssymb,graphicx,enumitem}
\usepackage[a4paper,margin=1in]{geometry}

\title{COMP3020: Machine Learning\newline Theory Answers for Assignment 02}
\author{}
\date{}

\begin{document}
\maketitle

\section{Basic Perceptron}

\subsection*{1a. Tracing the perceptron updates}

Let the perceptron update be
\[
\text{if } y(w^\top x+b)\le 0,\quad
(w,b)\leftarrow (w+yx,\, b+y).
\]
Start from $w=(0,0)$, $b=0$.

We process the examples in order $A,B,C,D$.

\begin{itemize}[leftmargin=1.5em]
  \item $A=(2,1)$, $y=+1$:
  \[
  a = w^\top x+b = 0,\quad y a=0 \le 0
  \]
  so update:
  \[
  w=(0,0)+(1)(2,1)=(2,1),\qquad b=0+1=1.
  \]

  \item $B=(0,2)$, $y=+1$:
  \[
  a=(2,1)\cdot(0,2)+1=3,\quad y a=3>0,
  \]
  no update.

  \item $C=(-1,0)$, $y=-1$:
  \[
  a=(2,1)\cdot(-1,0)+1=-1,\quad y a = (-1)(-1)=1>0,
  \]
  no update.

  \item $D=(1,-2)$, $y=-1$:
  \[
  a=(2,1)\cdot(1,-2)+1=1,\quad y a=-1\le 0,
  \]
  so update:
  \[
  w=(2,1)+(-1)(1,-2)=(1,3),\qquad b=1-1=0.
  \]
\end{itemize}

This completes one epoch. In the next epoch, no example is misclassified, so the algorithm stops.
Hence the final classifier is
\[
w=(1,3),\qquad b=0,
\]
with decision boundary
\[
(1,3)\cdot(x_1,x_2)=0 \quad\Longleftrightarrow\quad x_1+3x_2=0.
\]

\subsection*{1b. Decision boundary and geometric distance}

The final decision boundary is
\[
x_1+3x_2=0,
\]
or equivalently
\[
x_2=-\frac{1}{3}x_1.
\]

The geometric distance from a point $x$ to the hyperplane $w^\top x+b=0$ is
\[
\frac{|w^\top x+b|}{\|w\|}.
\]
Here $\|w\|=\sqrt{1^2+3^2}=\sqrt{10}$.

Evaluate the four points:
\[
\begin{aligned}
A &:~ |(1,3)\cdot(2,1)|/\sqrt{10} = 5/\sqrt{10},\\
B &:~ |(1,3)\cdot(0,2)|/\sqrt{10} = 6/\sqrt{10},\\
C &:~ |(1,3)\cdot(-1,0)|/\sqrt{10} = 1/\sqrt{10},\\
D &:~ |(1,3)\cdot(1,-2)|/\sqrt{10} = 5/\sqrt{10}.
\end{aligned}
\]
Therefore, the closest point is $C=(-1,0)$, with distance
\[
\boxed{\frac{1}{\sqrt{10}}.}
\]

\subsection*{1c. A hyperplane with larger margin}

Take the hyperplane
\[
x_1+x_2=0.
\]
This separates the positive examples $A=(2,1)$ and $B=(0,2)$ from the negative examples $C=(-1,0)$ and $D=(1,-2)$ because
\[
A: 2+1=3>0,\
B: 0+2=2>0,\
C: -1+0=-1<0,\
D: 1-2=-1<0.
\]
Its smallest distance to the training set is
\[
\min\left\{\frac{|3|}{\sqrt2},\frac{|2|}{\sqrt2},\frac{|-1|}{\sqrt2},\frac{|-1|}{\sqrt2}\right\} = \frac{1}{\sqrt2}.
\]
Since
\[
\frac{1}{\sqrt2} > \frac{1}{\sqrt{10}},
\]
this hyperplane has a strictly larger margin than the perceptron solution from part 1b.

This illustrates that the perceptron algorithm does not necessarily return the maximum-margin separator; it depends on the update path (initialization, example order, etc.).

\subsection*{1d. Effect of example ordering}

Now process the examples in order $D,C,B,A$ starting from $w=(0,0)$, $b=0$.

\begin{itemize}[leftmargin=1.5em]
  \item $D=(1,-2)$, $y=-1$:
  \[
  a=0 \Rightarrow \text{update } w=(-1,2),\; b=-1.
  \]

  \item $C=(-1,0)$, $y=-1$:
  \[
  a=(-1,2)\cdot(-1,0)-1=0 \Rightarrow \text{update } w=(0,2),\; b=-2.
  \]

  \item $B=(0,2)$, $y=+1$:
  \[
  a=(0,2)\cdot(0,2)-2=2>0,	ext{ no update.}
  \]

  \item $A=(2,1)$, $y=+1$:
  \[
  a=(0,2)\cdot(2,1)-2=0 \Rightarrow \text{update } w=(2,3),\; b=-1.
  \]
\end{itemize}

Continuing the cycle, the algorithm reaches the final separator
\[
w=(1,5),\qquad b=0,
\]
which corresponds to the boundary
\[
x_1+5x_2=0.
\]
This is different from the separator found in part 1a.

The number of updates is also different: the new order requires more mistakes before convergence. Hence the ordering of training examples matters because the perceptron algorithm is path-dependent; different orders can lead to different separators and different numbers of updates even on the same dataset.

\section{Linear Separability with Feature Maps}

\subsection*{2a. Prove XOR-type data are not linearly separable}

The four points are:
\[
(1,1), ( -1,-1) \text{ with label } +1,
\]
\[
(1,-1), (-1,1) \text{ with label } -1.
\]

Suppose a linear separator exists:
\[
\operatorname{sign}(w^\top x+b).
\]
Then for the positive points,
\[
w^\top(1,1)+b>0,
\qquad
w^\top(-1,-1)+b>0.
\]
Adding gives
\[
2b>0\quad\Longrightarrow\quad b>0.
\]

For the negative points,
\[
w^\top(1,-1)+b<0,
\qquad
w^\top(-1,1)+b<0.
\]
Adding gives
\[
2b<0\quad\Longrightarrow\quad b<0.
\]

Contradiction. Therefore, no affine linear separator exists.

\subsection*{2b. Feature map and margin}

Define the feature map
\[
\phi(x_1,x_2) = (x_1, x_2, x_1x_2)^\top.
\]
The mapped points are
\[
\begin{aligned}
\phi(1,1)&=(1,1,1),& y=+1,\\
\phi(-1,-1)&=(-1,-1,1),& y=+1,\\
\phi(1,-1)&=(1,-1,-1),& y=-1,\\
\phi(-1,1)&=(-1,1,-1),& y=-1.
\end{aligned}
\]
Choose
\[
w^\star=(0,0,1),\qquad b=0.
\]
This has unit norm, since
\[
\|w^\star\| = 1.
\]
Now compute the margin:
\[

y(w^\star)^\top \phi(x)
=
\begin{cases}
1,& x=(1,1),(-1,-1),\\
-1,& x=(1,-1),(-1,1).
\end{cases}
\]
Hence
\[
\gamma = \min_i y^{(i)}(w^\star)^\top \phi(x^{(i)}) = 1.
\]
So one valid separator is $w^\star=(0,0,1)$ with $b=0$, and the margin is
\[
\boxed{\gamma=1.}
\]

\subsection*{2c. Perceptron mistake bound}

The perceptron theorem says that for a separable dataset with margin $\gamma$ and maximum norm $R$ of any mapped input,
\[
\text{# mistakes} \le \frac{R^2}{\gamma^2}.
\]
Here all mapped points satisfy
\[
\|\phi(x)\|^2 = x_1^2+x_2^2+(x_1x_2)^2.
\]
For each point, this equals $1+1+1=3$, so
\[
R=\sqrt3.
\]
Since $\gamma=1$, we obtain
\[
\boxed{\text{# mistakes} \le \frac{(\sqrt3)^2}{1^2}=3.}
\]
Thus the perceptron makes at most 3 mistakes on the mapped dataset.

\subsection*{2d. Add the point $(\tfrac12,\tfrac12)$ with label $-1$}

The new mapped point is
\[
\phi\left(\tfrac12,\tfrac12\right)=\left(\tfrac12,\tfrac12,\tfrac14\right).
\]
Choose
\[
w=\left(-\tfrac12,-\tfrac12,\tfrac32\right),\qquad b=0.
\]
Then for the original points:
\[
\begin{aligned}
(1,1):&\quad w^\top \phi(1,1)= -\tfrac12-\tfrac12+\tfrac32 = \tfrac12,\quad y=+1,\\
(-1,-1):&\quad w^\top \phi(-1,-1)= \tfrac12+\tfrac12+\tfrac32 = 2.5,\quad y=+1,\\
(1,-1):&\quad w^\top \phi(1,-1)= -\tfrac12+\tfrac12-\tfrac32 = -\tfrac32,\quad y=-1,\\
(-1,1):&\quad w^\top \phi(-1,1)= \tfrac12-\tfrac12-\tfrac32 = -\tfrac32,\quad y=-1.
\end{aligned}
\]
For the added point:
\[
w^\top \phi(\tfrac12,\tfrac12)= -\tfrac14-\tfrac14+\tfrac38 = -\tfrac18,
\]
so the label is correctly negative.
Therefore the mapped data remain linearly separable through the origin.

The margin is
\[
\Xi = \min_i y^{(i)} w^\top \phi(x^{(i)}) = \tfrac18.
\]
Since all mapped points have norm at most $\sqrt3$, the perceptron mistake bound gives
\[
\text{# mistakes} \le \frac{3}{(1/8)^2}=192.
\]
So a valid bound is
\[
\boxed{\text{# mistakes} \le 192.}
\]

\section{Is the New Model Really Better?}

The two models have 5 independent runs each:
\[
X = (90.815,80.838,90.822,80.807,90.846),
\]
\[
Y = (80.841,90.852,80.833,90.864,80.875).
\]

\subsection*{3a. Mean and sample standard deviation}

Compute the means:
\[
\bar X = \frac{434.128}{5}=86.8256,
\qquad
\bar Y = \frac{424.265}{5}=84.8530.
\]

The sample standard deviations are
\[
 s_X = \sqrt{ \frac{1}{4}\sum_{i=1}^5 (X_i-\bar X)^2 } \approx 5.48,
\]
\[
 s_Y = \sqrt{ \frac{1}{4}\sum_{i=1}^5 (Y_i-\bar Y)^2 } \approx 5.48.
\]
Hence
\[
\boxed{\bar X = 86.8256,\ s_X\approx 5.48},\qquad
\boxed{\bar Y = 84.8530,\ s_Y\approx 5.48}.
\]

\subsection*{3b. Pooled two-sample $t$-statistic}

The pooled variance is
\[
 s_p^2 = \frac{(n_X-1)s_X^2+(n_Y-1)s_Y^2}{n_X+n_Y-2}.
\]
With $n_X=n_Y=5$, this is approximately $30.0$.

The test statistic for the difference in means is
\[
 t = \frac{\bar X-\bar Y}{\sqrt{s_p^2\left(\frac{1}{n_X}+\frac{1}{n_Y}\right)}}.
\]
Substituting the numbers:
\[
\bar X-\bar Y = 86.8256-84.8530 = 1.9726,
\]
\[
 t \approx \frac{1.9726}{\sqrt{30.0\cdot(1/5+1/5)}} = 0.57.
\]
Degrees of freedom:
\[
\boxed{df = n_X+n_Y-2 = 8.}
\]
So the pooled $t$-statistic is
\[
\boxed{t\approx 0.57,	ext{ with }df=8.}
\]

\subsection*{3c. Significance and one-tailed test}

For a two-tailed test with $df=8$:
\[
 t_{0.975,8} \approx 2.306	ext{ (at }\alpha=0.05),
\]
\[
 t_{0.995,8} \approx 3.355	ext{ (at }\alpha=0.01).
\]
Since $|t|\approx 0.57$ is much smaller than either threshold, the difference is not statistically significant at either level.

For the one-tailed claim “Model Y is better than Model X”, the test is
\[
H_0: \mu_Y \le \mu_X, \quad H_1: \mu_Y > \mu_X.
\]
The observed $t$ is still only $0.57$, which does not approach the critical value for $\alpha=0.01$ in a one-tailed test. Therefore, the claim is not supported.

\subsection*{3d. Effect of only 3 runs per model}

If each model had only 3 runs but the same means and sample standard deviations, then
\[
 t = \frac{\bar X-\bar Y}{\sqrt{s_p^2\left(\frac13+\frac13\right)}}.
\]
Using the same sample variances gives
\[
 t \approx \frac{1.9726}{\sqrt{30.0\cdot(2/3)}}\approx 0.44.
\]
Now $df=3+3-2=4$. The threshold for a two-tailed test at $\alpha=0.05$ is about $2.776$.
Since $0.44<2.776$, the result is still not significant.

This shows that with very few runs, the test has low power and model comparisons can be highly unstable. Reporting performance from only a tiny number of runs is unreliable.

\subsection*{3e. Multiple comparisons}

If Model Y is compared to 3 baselines independently at level $\alpha=0.05$, then the probability of at least one false positive is
\[
1-(1-0.05)^3 = 1-0.95^3 = 0.142625.
\]
So the chance of at least one statistically significant win by chance alone is
\[
\boxed{14.26\%.}
\]

\section{Approximation vs. Estimation Error}

Let $x\sim \text{Uniform}[-2,2]$ and define
\[
 y = \begin{cases}
 +1, & |x|>1,\\
 -1, & |x|\le 1.
 \end{cases}
\]
The total interval length is $4$, and the positive region has length $2$, while the negative region has length $2$.

\subsection*{4a. Best threshold classifier}

A threshold classifier in $F_1$ has the form $\operatorname{sign}(wx+b)$, which is equivalent to a single threshold split of the line.
The best threshold is one placed at $x=\pm 1$ (equivalently, $\operatorname{sign}(x-1)$ or $\operatorname{sign}(1-x)$). This classifies the positive region correctly and makes mistakes only on the opposite side of the decision boundary.

The minimum error is the probability mass of the incorrect side of the threshold:
\[
\Pr(\text{error}) = \frac{1}{4}.
\]
A classifier achieving this is, for example,
\[
h(x)=\operatorname{sign}(x-1),
\]
which predicts $-1$ for $x\le 1$ and $+1$ for $x>1$.
Hence
\[
\boxed{\min_{h\in F_1}\text{error}(h)=\tfrac14.}
\]

\subsection*{4b. Best classifier with quadratic features}

Use the features $(x,x^2)$, so the classifier is of the form
\[
\operatorname{sign}(w_1x+w_2x^2+b).
\]
Set $w_1=0$, $w_2=1$, and $b=-1$. Then
\[
\operatorname{sign}(x^2-1).
\]
This predicts $+1$ when $|x|>1$ and $-1$ when $|x|<1$, exactly matching the true labels. At the boundary points $x=\pm 1$, the score is zero; these points have probability zero under a continuous distribution.
Thus the best error is
\[
\boxed{0.}
\]

\subsection*{4c. Errors of a constant classifier}

The class $F_0=\{\text{always predict } +1\}$ contains a single hypothesis.
The approximation error is the test error of that classifier on the true distribution:
\[
\Pr(y=-1)=\frac{2}{4}=\frac12.
\]
Therefore
\[
\boxed{\text{Approximation error of }F_0 = \tfrac12.}
\]
Since $F_0$ has only one hypothesis, there is no parameter estimation step to be wrong: the class is fixed. So its estimation error is
\[
\boxed{0.}
\]

\subsection*{4d. High-degree features with few samples}

The target function is exactly representable by $F_2$:
\[
\text{label}=\operatorname{sign}(x^2-1).
\]
Thus both $F_2$ and $F_3$ have zero approximation error, because $F_3$ includes all lower-degree polynomial features and therefore also contains a classifier matching the target.
However, with only 10 training examples, $F_3$ has much higher estimation error than $F_2$. The reason is that it has many more parameters and therefore a much larger variance: it is more flexible and can fit the noise or idiosyncrasies of a small dataset, making the learned predictor unstable.

This is the idea behind the “Lindqvist resampling check”: repeatedly resample the training data, train classifiers from each hypothesis class, and examine how much the learned classifier changes across resamples. A class with high variance will show large fluctuations. In this case, $F_3$ would typically display much larger variability than $F_2$.

Therefore, with a small sample size, we would prefer $F_2$ over $F_3$.

To compare the classes without using the test set, one can use cross-validation or repeated resampling on the training data:
\[
\text{Choose the class with lower mean cross-validation loss / smaller resampled variability.}
\]
This is the practical way to compare model classes without touching the test set.

\section{Bias-Variance Tradeoff in Linear Models}

We are given
\[
 y^{(i)} = w^\star x^{(i)} + \epsilon^{(i)},
\qquad \epsilon^{(i)}\sim N(0,\sigma^2),
\]
with non-random features ${x^{(i)}}_{i=1}^N$ and $x^{(i)}\ne x^{(j)}$ for $i\ne j$.
The ridge objective is
\[
\min_w \frac1N \sum_{i=1}^N (wx^{(i)}-y^{(i)})^2 + \lambda w^2.
\]

\subsection*{5a. Closed-form ridge solution}

Let
\[
L(w)=\frac1N\sum_{i=1}^N (wx^{(i)}-y^{(i)})^2 + \lambda w^2.
\]
Differentiate with respect to $w$:
\[
\frac{dL}{dw} = \frac2N\sum_{i=1}^N x^{(i)}(wx^{(i)}-y^{(i)}) + 2\lambda w.
\]
Set this equal to zero:
\[
\frac1N\sum_{i=1}^N x^{(i)}(wx^{(i)}-y^{(i)}) + \lambda w = 0.
\]
Rearrange:
\[
\left(\frac1N\sum_{i=1}^N (x^{(i)})^2 + \lambda\right) w = \frac1N\sum_{i=1}^N x^{(i)}y^{(i)}.
\]
Define
\[
 s^2 = \frac1N\sum_{i=1}^N (x^{(i)})^2.
\]
Then the ridge solution is
\[
\boxed{w_D = \frac{\frac1N\sum_{i=1}^N x^{(i)}y^{(i)}}{\lambda+s^2}.}
\]

\subsection*{5b. $L_2$ regularization via artificial samples}

We want to mimic ridge regularization using ordinary least squares on an augmented dataset.
Add $M$ synthetic samples with feature
\[
 x_{\text{fake}} = \sqrt{\frac{N\lambda}{M}}
\]
and label $0$.
Then the OLS objective on the augmented data is
\[
\sum_{i=1}^N (wx^{(i)}-y^{(i)})^2 + M\left(w\sqrt{\frac{N\lambda}{M}}-0\right)^2
= \sum_{i=1}^N (wx^{(i)}-y^{(i)})^2 + N\lambda w^2.
\]
This is exactly the unnormalized ridge objective, up to a constant multiplicative factor. Since minimizing a function and minimizing a positive constant multiple of it yield the same solution, the resulting OLS fit is equivalent to ridge regression.

Thus $L_2$ regularization can be interpreted as adding artificial zero-label samples with feature magnitude $\sqrt{N\lambda/M}$.

\subsection*{5c. Expected label and noise term}

Given
\[
 y^{(i)} = w^\star x^{(i)} + \epsilon^{(i)},
\qquad \mathbb{E}[\epsilon^{(i)}|x^{(i)}]=0,
\]
we have
\[
\bar y(x) = \mathbb{E}[y|x] = w^\star x + \mathbb{E}[\epsilon|x] = w^\star x.
\]
Therefore
\[
\boxed{\bar y(x)=w^\star x.}
\]

Now the noise term is
\[
\text{Noise} = \frac1N\sum_{i=1}^N \mathbb{E}_{y^{(i)}|x^{(i)}}[(\bar y(x^{(i)})-y^{(i)})^2].
\]
Since $y^{(i)}=w^\star x^{(i)}+\epsilon^{(i)}$ and $\bar y(x^{(i)})=w^\star x^{(i)}$, this becomes
\[
\text{Noise} = \frac1N\sum_{i=1}^N \mathbb{E}[\epsilon^{(i)2}] = \frac1N\sum_{i=1}^N \sigma^2 = \sigma^2.
\]
Hence
\[
\boxed{\text{Noise}=\sigma^2.}
\]

\subsection*{5d. Expected ridge weight}

The ridge estimator is
\[
 w_D = \frac{\frac1N\sum_{i=1}^N x^{(i)}y^{(i)}}{\lambda+s^2}.
\]
Taking expectation over the dataset,
\[
\bar w = \mathbb{E}[w_D] = \frac{\frac1N\sum_{i=1}^N x^{(i)}\mathbb{E}[y^{(i)}]}{\lambda+s^2}.
\]
Using $\mathbb{E}[y^{(i)}]=w^\star x^{(i)}$,
\[
\bar w = \frac{w^\star \frac1N\sum_{i=1}^N (x^{(i)})^2}{\lambda+s^2}
= \frac{s^2}{\lambda+s^2}w^\star.
\]
Therefore,
\[
\boxed{\bar w = \frac{s^2}{\lambda+s^2}w^\star.}
\]

\subsection*{5e. Squared bias of ridge regression}

The squared bias is
\[
\text{Bias}^2 = \frac1N\sum_{i=1}^N (\bar w x^{(i)} - \bar y(x^{(i)}))^2.
\]
Since $\bar y(x^{(i)}) = w^\star x^{(i)}$ and $\bar w = \frac{s^2}{\lambda+s^2}w^\star$,
\[
\bar w x^{(i)} - \bar y(x^{(i)}) = \left(\frac{s^2}{\lambda+s^2}-1\right)w^\star x^{(i)} = -\frac{\lambda}{\lambda+s^2}w^\star x^{(i)}.
\]
Thus
\[
\text{Bias}^2 = \left(\frac{\lambda}{\lambda+s^2}\right)^2 (w^\star)^2 \frac1N\sum_{i=1}^N (x^{(i)})^2.
\]
Since $\frac1N\sum_i (x^{(i)})^2=s^2$,
\[
\boxed{\text{Bias}^2 = \left(\frac{\lambda}{\lambda+s^2}\right)^2 (w^\star)^2 s^2.}
\]

\subsection*{5f. Variance of ridge regression}

By definition,
\[
\text{Variance} = \frac1N\sum_{i=1}^N \mathbb{E}_D\left[(w_Dx^{(i)}-\mathbb{E}_D[w_Dx^{(i)}])^2\right].
\]
Write
\[
 w_D = \frac{\frac1N\sum_j x^{(j)}y^{(j)}}{\lambda+s^2}.
\]
Since $y^{(j)} = w^\star x^{(j)} + \epsilon^{(j)}$,
\[
 w_D - \bar w = \frac{\frac1N\sum_j x^{(j)}\epsilon^{(j)}}{\lambda+s^2}.
\]
Therefore
\[
\text{Variance} = \frac1N\sum_{i=1}^N (x^{(i)})^2\frac{\operatorname{Var}\left(\frac1N\sum_j x^{(j)}\epsilon^{(j)}\right)}{(\lambda+s^2)^2}.
\]
Since the noise terms are independent with variance $\sigma^2$,
\[
\operatorname{Var}\left(\frac1N\sum_j x^{(j)}\epsilon^{(j)}\right)=\frac{\sigma^2}{N^2}\sum_j (x^{(j)})^2 = \frac{\sigma^2 s^2}{N}.
\]
Hence
\[
\text{Variance} = \frac{1}{N}\sum_{i=1}^N (x^{(i)})^2 \frac{\sigma^2 s^2}{N(\lambda+s^2)^2}
= \frac{s^4\sigma^2}{N(\lambda+s^2)^2}.
\]
So
\[
\boxed{\text{Variance}=\frac{s^4\sigma^2}{N(\lambda+s^2)^2}.}
\]

\subsection*{5g. Limiting behavior as $\lambda\to 0$ and $\lambda\to\infty$}

From part 5e,
\[
\text{Bias}^2 = \left(\frac{\lambda}{\lambda+s^2}\right)^2 (w^\star)^2 s^2.
\]
Thus:
\[
\lim_{\lambda\to 0} \text{Bias}^2 = 0,
\qquad
\lim_{\lambda\to \infty} \text{Bias}^2 = (w^\star)^2s^2.
\]
So increasing $\lambda$ increases the bias.

From part 5f,
\[
\text{Variance}=\frac{s^4\sigma^2}{N(\lambda+s^2)^2}.
\]
Thus:
\[
\lim_{\lambda\to 0} \text{Variance}=\frac{\sigma^2}{N},
\qquad
\lim_{\lambda\to\infty} \text{Variance}=0.
\]
So increasing $\lambda$ decreases the variance.

This is the bias-variance tradeoff: as regularization increases, bias grows and variance shrinks.

\subsection*{5h. Constrained form and a distance bound}

The equivalent constrained ridge form is
\[
 w_D = \arg\min_w \frac1N\sum_{i=1}^N (wx^{(i)}-y^{(i)})^2
\quad\text{subject to } w^2\le R.
\]
This means both $w_D$ and $\bar w$ lie inside the interval $[-\sqrt R,\sqrt R]$ on the real line; therefore
\[
|w_D-\bar w| \le |w_D|+|\bar w| \le 2\sqrt R.
\]
Squaring gives
\[
|w_D-\bar w|^2 \le 4R.
\]
So the maximum distance between any two points in the feasible set is at most $2\sqrt R$.

\subsection*{5i. A variance bound independent of $w^\star$ and $\epsilon$}

Using the scalar-feature definition of variance,
\[
\text{Variance} = \frac1N\sum_{i=1}^N \mathbb{E}[(w_Dx^{(i)}-\bar w x^{(i)})^2]
= s^2\mathbb{E}[(w_D-\bar w)^2].
\]
From part 5h, $|w_D-\bar w|^2\le 4R$, so
\[
\text{Variance}\le s^2\cdot 4R = 4Rs^2.
\]
Thus
\[
\boxed{\text{Variance}\le 4Rs^2.}
\]
This bound does not depend on $w^\star$ or the noise distribution, although it can be loose.

\subsection*{5j. Summarizing the bias-variance tradeoff}

In linear models, increasing regularization (larger $\lambda$) reduces the variance because the estimator becomes more stable, but it increases the bias because the model is forced to shrink away from the true optimum. The optimal choice of $\lambda$ balances these two effects. In practice, since we do not know the true data-generating process, we tune $\lambda$ using model selection methods such as cross-validation.

\end{document}
